%

%

%%%%

%%%   Самарский государственный технический университет

%%%   Кафедра прикладной математики и информатики

%%%

%%%   Преамбула для статьи в журнал "Вестник Самарского государственного

%%%   технического университета. Серия: «Физико-математические науки»"

%%%   Ориентирована под MikTEx 2.4 for Win32

%%%

%%%   Хотя сам журнал собирается под GNU/Linux на дистрибутиве TeXLive



\documentclass[11pt,a4paper,twoside]{article}



\usepackage[cp1251]{inputenc}

\usepackage[T2A]{fontenc}

\usepackage[english,russian]{babel}

\usepackage{samgtubib}

\usepackage[dvips]{graphicx}

\usepackage[multiple]{footmisc}



\usepackage{samgtu}

\renewcommand{\VestnikYear}{2009} % Год издания Вестника

\renewcommand{\VestnikNumber}{2\,(19)} % Номер Вестника

\renewcommand{\VestnikMonth}{Ноябрь} % Месяц выхода Вестника

\renewcommand{\VestnikData}{01.11.09} % Дата подписания Вестника в печать

\renewcommand{\VestnikUPL}{26,64} % Усл. печ. л.

\renewcommand{\VestnikUIL}{25,73} % Уч.-изд. л.

\renewcommand{\VestnikRegNumber}{479/09} % Регистрационный номер

\renewcommand{\VestnikOrder}{994} % Номер заказа







%\usepackage{pst-barcode}

%\usepackage{auto-pst-pdf}

%

%\newcommand{\totop}[1]{\raisebox{6pt}[1pt][2pt]{#1}} 

%\newcommand{\tottop}[1]{\raisebox{12pt}[1pt][2pt]{#1}} 

%\newcommand{\totttop}[1]{\raisebox{18pt}[1pt][2pt]{#1}} 

%\newcommand{\tottttop}[1]{\raisebox{24pt}[1pt][2pt]{#1}} 

%\newcommand{\totttttop}[1]{\raisebox{30pt}[1pt][2pt]{#1}} 

%\newcommand{\tobot}[1]{\raisebox{-6pt}[1pt][2pt]{#1}} 

%\newcommand{\tobbot}[1]{\raisebox{-12pt}[1pt][2pt]{#1}} 

%\newcommand{\tobbbot}[1]{\raisebox{-18pt}[1pt][2pt]{#1}}

%

%%\samgtupubl % для генерации pdf, отдаваемого в типографию СамГТУ

%\pdfcrop % обрезка pdf для размещения на math-net.ru

%\draft % На корректуре

%%\nofilllastpage % Последняя русскоязычная страница заполнена не полностью

%\filllastpage

%%\briefcomm % Статья является кратким сообщением



\vsgtu{1328}



\FirstPage{7}

\LastPage{20}

%

%

%

%

%\begin{document}







%\hfill \begin{pspicture}(1in,1in)

%\psbarcode{http://dx.doi.org/10.14498/vsgtu1328}{}{qrcode}

%\end{pspicture}

%\vspace{-1in}



%\Partition{Механика деформируемого твердого тела}{}



\udk{517.956.6}

\msc{35M10, 35M12}



\rutitle{Нелокальная задача для уравнения смешанного~типа с~сингулярным коэффициентом~в~области, 

гиперболическая~часть~которой~--- вертикальная~полуполоса}

{Нелокальная задача для уравнения смешанного типа с~сингулярным коэффициентом  \dots}



\ruauthor{А.~А.~Абашкин}{А\,б\,а\,ш\,к\,и\,н~А.~А.}





\enauthor{A.~A.~Abashkin}{A\,b\,a\,s\,h\,k\,i\,n~A.~A.}





\ruaffil{\rusgasu.}

\enaffil{\ensgasu.}





\ruauthordetails{1}{\fullauthor{\rupid{60545}{Антон Александрович Абашкин}} \regalia{к.ф.-м.н.;  \mem{samcocaa@rambler.ru}}\position{старший преподаватель} \dept{каф. высшей математики}}





\enauthordetails{1}{\fullauthor{\enpid{60545}{Anton~A.~Abashkin}}

\regalia{Cand. Phys. \& Math. Sci.; \mem{samcocaa@rambler.ru}}\position{Assistant} \dept{Dept. of High Mathematics}}





\rukeywords{уравнение смешанного типа, уравнение с сингулярным коэффициентом, биортогональный ряд, базис Рисса, функции Бесселя}



\ruabstract{Для уравнения смешанного типа с сингулярным коэффициентом при младшей производной и спектральным параметром в области, гиперболическая часть которой "--- вертикальная полуполоса, а эллиптическая "--- прямоугольник, поставлена нелокальная задача с условием, связывающим значения искомой функции на правой и левой границах полуполосы и прямоугольника. 

При этом на линии изменения типа от искомой функции требуется лишь непрерывность.  

Для исследования поставленной задачи применен спектральный метод. Доказаны теоремы единственности и существования решения исследуемой задачи. 

Решение построено в виде разложения в биортогональный ряд по одной системе тригонометрических функций, предложенной в работах Е.~И.~Моисеева при этом, коэффициенты разложения получены как решения соответствующих систем ОДУ. 

Дано обоснование равномерной сходимости соответствующих рядов при определенных ограничениях на условия задачи.

}



\entitle{A Nonlocal Problem for Mixed Type Equation with~Singular Coefficient in Domain with~Half-Strip~as~Hyperbolic~Part}







\enabstract{A nonlocal problem for mixed type equation with a singular coefficient and the spectral parameter is formulated in the field, which hyperbolic part is vertical half-strip and elliptic part is rectangle. The nonlocal condition of problem combines the values of required function on the right and left boundaries of half-stripe and rectangle.

The only requirement on the unknown function in the change type line is continuity. To research the given problem we apply the spectral method. The uniqueness and existence of a solution are proved. The solution is constructed as biortogonal series. Coefficients of this series should require special ODE systems, solved in the paper. The uniform convergence of the series is proved with the restrictions on problem conditions.}



\enkeywords{mixed type equation, equation with a singular coefficient, biorthogonal series, Riesz basis, Bessel functions}



\date{14/VII/2014}

\revisiondate{16/VIII/2014}

\accepteddate{27/VIII/2014}



\makerutitle







\Section[n]{Постановка задачи}

Для уравнения 

\begin{equation}\label{ab:eq1}

Lu=u_{xx}+\mathop{\rm sign} y\cdot u_{yy}+\frac{2p}{|y|}u_y+k u=0, \quad p\ge\frac12,\quad k\in\mathbb{R}

\end{equation}

рассмотрим следующую задачу.



\smallskip



{\small \sc Задача.} \it В области 

$$D=\{(x,y)\mid 0<x<1,\; y<\alpha\},\quad  \alpha>0$$ найти функцию $u(x,y),$ удовлетворяющую следующим условиям{\/\rm:}

\begin{gather}

\label{ab:eq2}

u\in C(\overline{D})\cap C^2(D^+\cup D^-),\quad Lu=0;

\\

\label{ab:eq3}

 u(0, y)=u(1, y),\quad u_x(0, y)=0,\quad \text{при} \quad y<\alpha;

\\

\label{ab:eq4}

u(x,\alpha)=\varphi(x) \quad \text{при} \quad 0<x<1,

\end{gather}

где $D^{+}=D\cap\{y>0\},$ $D^{-}=D\cap\{y<0\};$ $\varphi(x)$ "--- заданная непрерывная функция$,$ удовлетворяющая условию $\varphi(0)=\varphi(1).$



\rm



\smallskip



Отметим, что уравнение \eqref{ab:eq1} в области эллиптичности представляет собой обобщенное осесимметрическое уравнение Гельмгольца, а в области гиперболичности содержит в себе (при $k=0$) уравнение Эйлера"--~Пуассона"--~Дарбу. 

Для частного случая уравнения \eqref{ab:eq1} при $k=0$ краевые задачи в~областях различного вида рассматривались в публикациях С.~П.~Пулькина \cite{ab:bib:Pulkin}, В.~Ф.~Волкодавова \cite{ab:bib:Volkodavov}, О.~А.~Репина \cite{ab:bib:Repin}, М.~Х.~Рузиева \cite{ab:bib:Ruziev} и др.



В работах \cite{ab:bib:Sabitov, ab:bib:SabitovSuleimanova} был предложен новый подход для доказательства единственности и~построения решений краевых задач для уравнений смешанного типа "--- спектральный метод. 

Ранее подобный метод был применен В.~А.~Ильиным и Е.~И.~Моисеевым~\cite{ab:bib:Il'inMoiseev} при рассмотрении нелокальной задачи с~условием, похожим на первое равенство условия~\eqref{ab:eq3}. 



Отметим, что впервые задача с нелокальным условием типа \eqref{ab:eq3} возникла в статье  Ф.~И.~Франкля~\cite{ab:bib:Francl} при изучении газодинамической задачи об обтекании профилей потоком дозвуковой скорости со сверхзвуковой зоной, оканчивающейся прямым скачком уплотнения. 

Существует довольно много работ, в~которых изучены краевые задачи с условием \eqref{ab:eq3} для различных уравнений, например [\citen{ab:bib:Moiseev2}--\citen{ab:bib:Sidorenko}], в частности, в публикации \cite{ab:bib:Sidorenko}, как и в данном исследовании, рассматривается нелокальная задача для уравнения смешанного типа в полуполосе. 

Задачи с~нелокальным условием \eqref{ab:eq3} для уравнения \eqref{ab:eq1} в области эллиптичности были исследованы в~работах М.~Е.~Лернера и О.~А.~Репина~\cite{ab:bib:LernerRepin1},  Е.~И.~Моисеева \cite{ab:bib:Moiseev1} и автора настоящей работы~\cite{ab:bib:Abashkin1,  ab:bib:Abashkin2}. 

В~работах [\citen{ab:bib:Moiseev1}--\citen{ab:bib:Abashkin2}] исследование основано на спектральном методе.



Для уравнения Лаврентьева"--~Бицадзе задача с~условием \eqref{ab:eq3} в прямоугольной области была рассмотрена Ю.~К.~Сабитовой в~работе \cite{ab:bib:Sabitova1}, а для уравнения смешанного типа $K(y)u_{xx}+u_{yy}-b^2K(y)u=0$  "--- в~\cite{ab:bib:Sabitova2}. 



В данной работе тем же методом, что и в статьях \cite{ab:bib:Moiseev2, ab:bib:Moiseev1, ab:bib:Sabitova1}, доказана единственность решения задачи \eqref{ab:eq2}--\eqref{ab:eq4}. 

Решение построено в виде суммы ряда. 

Существование решения задачи доказано путем установления равномерной сходимости соответствующих рядов.



Отметим, что в задаче \eqref{ab:eq2}--\eqref{ab:eq4} на линии  $y=0$ требуется лишь непрерывность искомой функции.



\smallskip



\Section{Единственность решения задачи}



\hypertarget{ab:th1}{}



{\small \sc Теорема 1.} \it Если решение задачи \eqref{ab:eq2}--\eqref{ab:eq4} существует$,$ то оно единственно\rm. 



\smallskip



\begin{proof}

Следуя работе \cite{ab:bib:Moiseev1}, рассмотрим функции

\begin{gather}

\label{ab:eq5}

v_n(y)=4\int _0^1u(x,y)\sin(2\pi nx)dx,

\\

\label{ab:eq6}

u_n(y)=4\int _0^1u(x,y)(1-x)\cos(2\pi nx)dx,\quad u_0(y)=2\int _0^1u(x,y)(1-x)dx,

\end{gather}

где $u(x,y)$ "--- решение задачи \eqref{ab:eq2}--\eqref{ab:eq4}.



Так как $u(x,y)$ является решением уравнения \eqref{ab:eq1}, то справедливо равенство

\begin{equation}

\label{ab:eq8}

4\int _0^1 \Bigl(u_{xx}+\mathop{\rm sign} y\cdot u_{yy}+\frac{2p}{|y|}u_y+k u\Bigr)\sin(2\pi nx)dx=0.

\end{equation}



Дважды проинтегрировав по частям, с учетом условия \eqref{ab:eq3} и определения \eqref{ab:eq5} соотношение \eqref{ab:eq8} запишем в виде

\begin{equation}

\label{ab:eq9}

4\int _0^1u_{xx}\sin(2\pi nx)dx=-(2\pi n)^2v_n(y).

\end{equation}



Из равенств \eqref{ab:eq8} и \eqref{ab:eq9} и определения \eqref{ab:eq5} следует

\begin{gather}

\label{ab:eq10}

v_n''+\frac{2p}{y}v_n'-\mathop{\rm sign} y\cdot\sigma_n^2v_n=0 \quad \text{при}\quad (2\pi n)^2>k;

\\

\label{ab:eq50}

v_n''+\frac{2p}{y}v_n'+\mathop{\rm sign} y\cdot\sigma_n^2v_n=0 \quad \text{при}\quad (2\pi n)^2<k;

\\

\label{ab:eq60}

v_n''+\frac{2p}{y}v_n'=0 \quad \text{при}\quad (2\pi n)^2=k,

\end{gather}

где $\sigma_n=\sqrt{|(2\pi n)^2-k|}$.



При $y>0$  уравнение \eqref{ab:eq10} заменой $v_n(y)=y^{-p_1}F(\sigma_n y)$,  $p_1=p-1/2$, сводится к модифицированному уравнению Бесселя \cite[c.~141]{ab:bib:Lebedev} и его общее решение представимо в виде

\begin{equation}\label{ab:eq11}

v_n=A_ny^{-p_1}I_{p_1}(\sigma_ny)+B_ny^{-p_1}K_{p_1}(\sigma_ny)\quad \text{при} \quad y>0,

\end{equation}

где $I_{\nu}(z)$ и $K_{\nu}(z)$ "--- модифицированные функции Бесселя \cite[c.~139]{ab:bib:Lebedev}.









Вследствие  асимптотики \cite[c.~173]{ab:bib:Lebedev}

\begin{equation}

\label{ab:eq46}

K_{\nu}(z)\simeq\frac{2^{\nu-1}\Gamma(\nu)}{z^{\nu}},\quad z\to0\quad \nu>0;

\qquad K_0(z)\simeq\ln\frac{2}{z},\quad z\to 0,

\end{equation}

для ограниченности функции $u(x,y)$ вблизи нуля необходимо, чтобы $B_n=0$.



Из соотношения \eqref{ab:eq4} получаем следующее условие:

\begin{equation}

\label{ab:eq12}

v_n(\alpha)=a_n,

\end{equation}

где $a_n= \displaystyle \int _0^1\varphi(x)\sin(2\pi nx)dx$.







В силу условия \eqref{ab:eq12} имеем

$$A_n=\frac{a_n\alpha^{p_1}}{I_{p_1}(\sigma_n\alpha)}.$$



Окончательно функция $v_n(y)$ при $y>0$ и $(2\pi n)^2>k$ принимает вид

\begin{equation}\label{ab:eq13}

v_n(y)=\frac{a_n\alpha^{p_1}}{I_{p_1}(\sigma\alpha)}y^{-p_1}I_{p_1}(\sigma_ny).

\end{equation}



При $y>0$  уравнение \eqref{ab:eq50} заменой $v_n(y)=y^{-p_1}F(\sigma_n y)$ сводится к уравнению Бесселя \cite[c.~134]{ab:bib:Lebedev} и его общее решение представимо в виде

\begin{equation}\label{ab:eq52}

v_n(y)=C_ny^{-p_1}J_{p_1}(\sigma_n y)+D_ny^{-p_1}Y_{p_1}(\sigma_n y),

\end{equation}

где $J_{\nu}(z)$ и $Y_{\nu}(z)$ "--- функции Бесселя первого \cite[c.~132]{ab:bib:Lebedev} и второго \cite[c.~134]{ab:bib:Lebedev} рода соответственно.





Из требования ограниченности функции $u(x,y)$ вблизи нуля и асимптотики \cite[c.~172]{ab:bib:Lebedev}

\begin{equation}\label{ab:eq47} Y_{\nu}(z)\simeq-\frac{2^{\nu}\Gamma(\nu)}{\pi z^{\nu}},\quad z\to0,\quad \nu>0;\qquad Y_0(z)\simeq-\frac2\pi\ln\frac2z,\quad z\to0,

\end{equation}

 следует, что $D_n=0$. 

 

Вследствие условия \eqref{ab:eq12} имеем $$C_n=\frac{a_n\alpha^{p_1}}{J_{p_1}(\sigma_n\alpha)}.

$$



Окончательно при $y>0$ и $(2\pi n)^2<k$ получаем

\begin{equation}

\label{ab:eq53}

v_n(y)=\frac{a_n\alpha^{p_1}}{J_{p_1}(\sigma\alpha)}y^{-p_1}J_{p_1}(\sigma_ny).

\end{equation}



При $y<0$  уравнение \eqref{ab:eq10} сводится к уравнению Бесселя и с~учетом требования ограниченности функции $u(x,u)$ вблизи нуля его решение имеет вид

\begin{equation}

\label{ab:eq54}

v_n(y)=E_n(-y)^{-p_1}J_{p_1}(-\sigma_n y).

\end{equation}

  

Используя непрерывность функции $u(x,y)$ и поведение функций $I_{\nu}(z)$, $J_{\nu}(z)$ при малых значениях $z$ \cite[c.~172, 173]{ab:bib:Lebedev}:

\begin{gather}

\label{ab:eq44}

I_{\nu}(z)\simeq\frac{z^{\nu}}{2^{\nu}\Gamma(1+\nu)},\quad z\to0,

\\

\label{ab:eq45}J_{\nu}(z)\simeq\frac{z^{\nu}}{2^{\nu}\Gamma(1+\nu)},\quad z\to0,

\end{gather}

 получаем $E_n=A_n$.



Окончательно $v_n(y)$ при $y<0$ и $(2\pi n)^2>k$ принимает вид

\begin{equation}

\label{ab:eq14}

v_n(y)=\frac{a_n\alpha^{p_1}}{I_{p_1}(\sigma\alpha)}(-y)^{-p_1}J_{p_1}(-\sigma_ny).

\end{equation}



Таким же образом получаем выражение для $v_n(y)$ при $y<0$ и $(2\pi n)^2<k$:

\begin{equation}\label{ab:eq61}

v_n(y)=\frac{a_n\alpha^{p_1}}{J_{p_1}(\sigma\alpha)}(-y)^{-p_1}I_{p_1}(-\sigma_ny).

\end{equation}



Уравнение \eqref{ab:eq60} имеет следующее общее решение:

$$v_n(y)=F_ny^{1-2p}+G_n\quad\mbox{при}\quad p>\frac12;\qquad v_n(y)=F_n\ln y+G_n\quad\mbox{при}\quad p=\frac12.

$$



С учетом требования ограниченности $u(x,y)$ вблизи нуля и условия \eqref{ab:eq4} при $(2\pi n)^2=k$ функция $v_n(y)$ будет иметь вид

\begin{equation}\label{ab:eq61}

v_n(y)=a_n.

\end{equation}



Аналогично получаем дифференциальные уравнения для $u_n(y)$:

\begin{gather}

\label{ab:eq15}

u_n''+\frac{2p}{y}u_n'-\mathop{\rm sign} y\cdot\sigma_n^2u_n=-\mathop{\rm sign} y\cdot4\pi nv_n\quad\mbox{при}\quad (2\pi n)^2>k;

\\

\label{ab:eq56}

u_n''+\frac{2p}{y}u_n'+\mathop{\rm sign} y\cdot\sigma_n^2u_n=-\mathop{\rm sign} y\cdot4\pi nv_n\quad\mbox{при}\quad (2\pi n)^2<k;

\\

\label{ab:eq62}

u_n''+\frac{2p}{y}u_n'=-\mathop{\rm sign} y\cdot4\pi nv_n\quad\mbox{при}\quad (2\pi n)^2=k.

\end{gather}





Уравнения \eqref{ab:eq15}--\eqref{ab:eq62} являются линейными неоднородным. 

Общее решение таких уравнений представимо в виде суммы общего решения однородного уравнения и частного решения неоднородного.



Однородное уравнение, соответствующее \eqref{ab:eq15}, совпадает с уравнением \eqref{ab:eq10}, поэтому его общее решение выражается формулами \eqref{ab:eq11} и \eqref{ab:eq14}.



Частным решением неоднородного уравнения при $y>0$ будет функция

$$

u_n(y)=-\frac{2\pi na_n\alpha^{p_1}}{\sigma_nI_{p_1}(\sigma_n\alpha)}y^{-p_1+1}I_{p_1-1}(\sigma_n y).

$$



Таким образом, общее решение уравнения \eqref{ab:eq15} при $y>0$ имеет вид

\begin{equation}\label{ab:eq17}

u_n(y)=F_ny^{-p_1}I_{p_1}(\sigma_n y)+G_ny^{-p_1}K_{p_1}(\sigma_n y)-\frac{2\pi na_n\alpha^{p_1}}{\sigma_nI_{p_1}(\sigma_n\alpha)}y^{-p_1+1}I_{p_1-1}(\sigma_n y).

\end{equation}



Ввиду асимптотики \eqref{ab:eq46} и требования непрерывности функции $u(x,y)$ необходимо, чтобы  $G_n=0$.



В силу условия \eqref{ab:eq4} имеем 

$$F_n=\frac{2\pi na_n\alpha^{p_1+1}I_{p_1-1}(\sigma_n\alpha)}{\sigma_nI_{p_1}^2(\sigma_n\alpha)}+\frac{b_n\alpha^{p_1}}{I_{p_1}(\sigma_n\alpha)},$$ 

где $b_n=4 \displaystyle \int _0^1\varphi(x)(1-x)\cos (2\pi nx)dx$.







Окончательно функция $u_n(y)$ при $y>0$ и $(2\pi n)^2>k$ принимает вид

 \begin{multline}\label{ab:eq18}

u_n(y)=\left(\frac{2\pi na_n\alpha^{p_1+1}I_{p_1-1}(\sigma_n\alpha)}{\sigma_nI_{p_1}^2(\sigma_n\alpha)}+\frac{b_n\alpha^{p_1}}{I_{p_1}(\sigma_n\alpha)}\right)y^{-p_1}I_{p_1}(\sigma_ny)-\\

-\frac{2\pi na_n\alpha^{p_1}}{\sigma_nI_{p_1}(\sigma_n\alpha)}y^{-p_1+1}I_{p_1-1}(\sigma_n y).

\end{multline} 



Частным решением уравнения \eqref{ab:eq15} при $y<0$ является функция

$$

u_n(y)=\frac{\pi na_n\alpha^{p_1}}{p_1\sigma_nI_{p_1}(\sigma_n\alpha)}(-y)^{-p_1+1}J_{p_1-1}(-\sigma_n y).

$$



Общим решением уравнения \eqref{ab:eq15} при $y<0$ является

$$

u_n(y)=T_n(-y)^{-p_1}J_{p_1}(-\sigma_n y)+\frac{\pi na_n\alpha^{p_1}}{p_1\sigma_nI_{p_1}(\sigma_n\alpha)}(-y)^{-p_1+1}J_{p_1-1}(-\sigma_n y).

$$



Приравнивая выражения для $u_n(y)$ при $y\to0+$ и $y\to0-$ и выражая коэффициент $T_n$ будем иметь 

$$T_n=\frac{2\pi na_n\alpha^{p_1+1}I_{p_1-1}(\sigma_n\alpha)}{\sigma_nI_{p_1}^2(\sigma_n\alpha)}+\frac{b_n\alpha^{p_1}}{I_{p_1}(\sigma_n\alpha)}-\frac{2p(2p_1+1)\pi na_n\alpha^{p_1}}{p_1\sigma_n^2I_{p_1}(\sigma_n\alpha)}.

$$



Окончательно $u_n(y)$ при $y<0$ и $(2\pi n)^2>k$ принимает вид

\begin{multline}

\label{ab:eq19}

u_n(y)=\Bigl(\frac{2\pi na_n\alpha^{p_1+1}I_{p_1-1}(\sigma_n\alpha)}{\sigma_nI_{p_1}^2(\sigma_n\alpha)}+\frac{b_n\alpha^{p_1}}{I_{p_1}(\sigma_n\alpha)}- \\

 - \frac{2p(2p_1+1)\pi na_n\alpha^{p_1}}{p_1\sigma_n^2I_{p_1}(\sigma_n\alpha)}\Bigr)(-y)^{-p_1}J_{p_1}(-\sigma_n y)+\\

+\frac{\pi na_n\alpha^{p_1}}{p_1\sigma_nI_{p_1}(\sigma_n\alpha)}(-y)^{-p_1+1}J_{p_1-1}(-\sigma_n y).

\end{multline}







Выполнив подобные вычисления применительно к уравнениям \eqref{ab:eq56} и \eqref{ab:eq62} получим

\begin{multline}

\label{ab:eq71}

u_n(y)=\Bigl(\frac{2\pi na_n\alpha^{p_1+1}J_{p_1-1}(\sigma_n\alpha)}{\sigma_nJ_{p_1}^2(\sigma_n\alpha)}+\frac{b_n\alpha^{p_1}}{J_{p_1}(\sigma_n\alpha)}\Bigr)y^{-p_1}J_{p_1}(\sigma_ny)-\\

-\frac{2\pi na_n\alpha^{p_1}}{\sigma_nJ_{p_1}(\sigma_n\alpha)}y^{-p_1+1}J_{p_1-1}(\sigma_n y)\quad\mbox{при}\quad y>0,\quad(2\pi n)^2<k;

\end{multline} 

\begin{multline}

\label{ab:eq71}

u_n(y)=\Bigl(\frac{2\pi na_n\alpha^{p_1+1}J_{p_1-1}(\sigma_n\alpha)}{\sigma_nJ_{p_1}^2(\sigma_n\alpha)}+\frac{b_n\alpha^{p_1}}{J_{p_1}(\sigma_n\alpha)}- \\ - \frac{2p(2p_1+1)\pi na_n\alpha^{p_1}}{p_1\sigma_n^2J_{p_1}(\sigma_n\alpha)}\Bigr)(-y)^{-p_1}I_{p_1}(-\sigma_n y)+\\

+\frac{\pi na_n\alpha^{p_1}}{p_1\sigma_nJ_{p_1}(\sigma_n\alpha)}(-y)^{-p_1+1}I_{p_1-1}(-\sigma_n y)\quad\mbox{при}\quad y<0,\quad(2\pi n)^2<k;

\end{multline}

\begin{equation}

\label{ab:eq71}

u_n(y)=b_n+\mathop{\rm sign} y\cdot\frac{\pi na_n}{p_1+1}\alpha^2-\mathop{\rm sign} y\cdot\frac{\pi na_n}{p_1+1}y^2 \quad\mbox{при}\quad(2\pi n)^2=k.

\end{equation}







Аналогичным образом (как и для $v_n(y)$ и $u_n(y)$) выражаем $u_0(y)$:

\begin{gather}

\label{ab:eq71}

u_0(y)=\frac{a_0\alpha^{p_1}}{I_{p_1}(k\alpha)}y^{-p_1}I_{p_1}(\sqrt{-k}y)\quad\mbox{при}\quad y>0,\quad k<0;

\\

\label{ab:eq71}

u_0=\frac{a_0\alpha^{p_1}}{I_{p_1}(k\alpha)}(-y)^{-p_1}J_{p_1}(-\sqrt{-k}y)\quad\mbox{при}\quad y<0,\quad k<0;

\\

\label{ab:eq71}

u_0(y)=\frac{a_0\alpha^{p_1}}{J_{p_1}(k\alpha)}y^{-p_1}J_{p_1}(\sqrt{k}y)\quad\mbox{при}\quad y>0,\quad k>0;

\\

\label{ab:eq71}

u_0=\frac{a_0\alpha^{p_1}}{J_{p_1}(k\alpha)}(-y)^{-p_1}I_{p_1}(-\sqrt{k}y)\quad\mbox{при}\quad y<0,\quad k>0;

\\

\label{ab:eq71}

u_0=a_0\quad\mbox{при}\quad k=0,

\end{gather}

где $a_0=2 \displaystyle \int _0^1\varphi(x)(1-x)dx$.



Из формул \eqref{ab:eq13}, \eqref{ab:eq53}, \eqref{ab:eq14}--\eqref{ab:eq61}, \eqref{ab:eq18}--\eqref{ab:eq71} следует, что если $\varphi(x)\equiv0$, то $u_n\equiv0$ и $v_n\equiv0$, но тогда $u(x,y)\equiv0$ в силу полноты системы функций \cite{ab:bib:Moiseev2}



$$\{4\sin 2\pi nx\}_{n=1}^{\infty},\quad \{2(1-x)\},\quad \{4(1-x)\cos 2\pi nx\}_{n=1}^{\infty},

$$

что завершает доказательство теоремы \hyperlink{ab:th1}{1}.

\end{proof}







\Section{Существование решения задачи}



\hypertarget{ab:th2}{}



{\small\sc Теорема 2.} \it Если $\varphi(x)\in C^2[0, 1]$ и $k\ne(2\pi n)^2+\left({r_m}/{\alpha}\right)^2,$ где $r_m$ "--- положительные нули функции $J_{p_1}(z),$ пронумерованные в порядке возрастания$,$ $n\in\mathbb{Z},$ то решение задачи \eqref{ab:eq2}--\eqref{ab:eq4} существует и представимо в виде суммы ряда

\begin{equation}

\label{ab:eq26}

u(x,y)=u_0(y)+\sum\limits_{n=1}^{\infty}u_n(y)\cos2\pi nx+\sum\limits_{n=1}^{\infty}v_n(y)x\sin2\pi nx,

\end{equation}

где функции $u_0(y),$ $u_n(y)$ и $v_n(y)$ определяются формулами \rm \eqref{ab:eq13}, \eqref{ab:eq53}, \eqref{ab:eq14}--\eqref{ab:eq61}, \eqref{ab:eq18}--\eqref{ab:eq71}.





\begin{proof}

Поскольку системы функций

$$

\{4\sin 2\pi nx\}_{n=1}^{\infty},\quad \{2(1-x)\},\quad \{4(1-x)\cos 2\pi nx\}_{n=1}^{\infty},

$$

$$

\{x\sin 2\pi nx\}_{n=1}^{\infty},\quad \{1\},\quad \{\cos 2\pi nx\}_{n=1}^{\infty},

$$

образуют базис Рисса в $L_{2}(0, 1)$ \cite{ab:bib:Moiseev2}, то достаточно доказать два положения:



\hypertarget{ab:1}{}

\begin{itemize}

\item [1)] равномерную сходимость на множестве $\overline{D}$ остатка ряда \eqref{ab:eq26}, содержащего члены с номерами, удовлетворяющими соотношению \hypertarget{ab:2}{} $(2\pi n)^2>k$; 



\item [2)] равномерную сходимость рядов, полученных из того же остатка ряда \eqref{ab:eq26} однократным и двукратным почленным дифференцированием по $x$ и по $y$ на множествах $D_{\varepsilon}=\{(x,y)\mid 0<x<1, \; y<\alpha-\varepsilon\}$, где $\varepsilon$ "--- сколь угодно малое положительное число.



\end{itemize}



Докажем выполнение положения \hyperlink{ab:1}{1)}.

Подставив вместо $u_n(y)$ и $v_n(y)$ их значения, определяемые формулами \eqref{ab:eq13}, \eqref{ab:eq14}, \eqref{ab:eq18} и \eqref{ab:eq19}, получим, что достаточно доказать равномерную сходимость следующих рядов:

\begin{gather}

s_1(x,y)=\sum\limits_{n=1}^{\infty}\frac{a_n\alpha^{p_1}}{I_{p_1}(\sigma\alpha)}y^{-p_1}I_{p_1}(\sigma_ny)x\sin2\pi nx,\quad y>0,

\\

s_1(x,y)=\sum\limits_{n=1}^{\infty}\frac{a_n\alpha^{p_1}}{I_{p_1}(\sigma\alpha)}(-y)^{-p_1}J_{p_1}(-\sigma_ny)x\sin2\pi nx,\quad y<0;

\\

s_2(x,y)=\sum\limits_{n=1}^{\infty}\Bigl(\frac{2\pi na_n\alpha^{p_1+1}I_{p_1-1}(\sigma_n\alpha)}{\sigma_nI_{p_1}^2(\sigma_n\alpha)}+\frac{b_n\alpha^{p_1}}{I_{p_1}(\sigma_n\alpha)}\Bigr) \times \\

\hspace{7cm} \times y^{-p_1}I_{p_1}(\sigma_ny)\cos2\pi nx,\quad y>0,

\\

s_2(x,y)=\sum\limits_{n=1}^{\infty}\Bigl(\frac{2\pi na_n\alpha^{p_1+1}I_{p_1-1}(\sigma_n\alpha)}{\sigma_nI_{p_1}^2(\sigma_n\alpha)}+\frac{b_n\alpha^{p_1}}{I_{p_1}(\sigma_n\alpha)}-\frac{2p(2p_1+1)\pi n\alpha^{p_1}}{p_1\sigma_n^2I_{p_1}(\sigma_n\alpha)}\Bigr)\times\\

\hspace{6cm}\times(-y)^{-p_1}J_{p_1}(-\sigma_ny)\cos2\pi nx,\quad y<0;

\\

s_3(x,y)=\sum\limits_{n=1}^{\infty}\frac{2\pi na_n\alpha^{p_1}}{\sigma_nI_{p_1}(\sigma_n\alpha)}y^{-p_1+1}I_{p_1-1}(\sigma_n y)\cos2\pi nx,\quad y>0,

\\

s_3(x,y)=\sum\limits_{n=1}^{\infty}\frac{\pi na_n\alpha^{p_1}}{p_1\sigma_nI_{p_1}(\sigma_n\alpha)}(-y)^{-p_1+1}J_{p_1-1}(-\sigma_n y)\cos2\pi nx,\quad y<0.

\end{gather}



Из формулы дифференцирования \cite[c.~141]{ab:bib:Lebedev} 

$$

\left(z^{-\nu}I_{\nu}(z)\right)'=z^{-\nu}I_{\nu+1}(z),

$$

асимптотики \eqref{ab:eq44} и отсутствия у функции $I_{\nu}(z)$ положительных нулей \cite[c.~173]{ab:bib:Lebedev} следует, что функция $z^{-\nu}I_{\nu}(z)$ положительна и монотонно возрастает при $z>0$, поэтому верны неравенства 

\begin{gather}

\label{ab:eq30}

\left|y^{-p_1}I_{p_1}(\sigma_ny)\right|\le\alpha^{-p_1}I_{p_1}(\sigma_n\alpha),\quad y>0;

\\

\label{ab:eq31}

\left|y^{-p_1-1}I_{p_1+1}(\sigma_ny)\right|\le\alpha^{-p_1-1}I_{p_1+1}(\sigma_n\alpha),\quad y>0.

\end{gather} 



Согласно асимптотикам \eqref{ab:eq45} и \cite[c.~172]{ab:bib:Lebedev} 

$$

J_{\nu}(z)\simeq\sqrt{\frac{2}{\pi z}}\cos\Bigl(z-\frac{\nu\pi}{2}-\frac{\pi}{4}\Bigr), \quad z\to+\infty,

$$

функция $z^{-\xi}J_{\nu}(z)$ при $\xi\le\nu$ ограничена при $z>0$.

Обозначим 

$$M_{J_{\nu}}=\max\limits_{z>0}z^{-\nu}J_{\nu}(z),$$ 

тогда справедливы оценки 

$$

\left|(-y)^{-p_1}J_{p_1}(-\sigma_ny)\right|<\sigma_n^{p_1}M_{J_{p_1}}, \quad \left|(-y)^{-p_1+1}J_{p_1-1}\right|<\sigma_n^{p_1-1}M_{J_{p_1-1}}.

$$ 







В силу асимптотики \cite[c.~173]{ab:bib:Lebedev}

\begin{equation}

\label{ab:eq28}

I_{\nu}(z)\simeq\frac{e^z}{\sqrt{2\pi z}},\quad z\to+\infty,

\end{equation}

существует такое $N$, что для $n>N$ выполняются свойства 

$$

\sigma_n^{p_1}M_{J_{p_1}}<\alpha^{-p_1}I_{p_1}(\sigma_n\alpha), \quad \sigma_n^{p_1-1}M_{J_{p_1-1}}<\alpha^{-p_1+1}I_{p_1-1}(\sigma_n\alpha)$$ и поэтому верны неравенства

\begin{gather}

\label{ab:eq29}

\left|(-y)^{-p_1}J_{p_1}(-\sigma_ny)\right|<\alpha^{-p_1}I_{p_1}(\sigma_n\alpha),\quad y<0,\quad n>N;

\\

\label{ab:eq32}

\left|(-y)^{-p_1+1}J_{p_1-1}(-\sigma_ny)\right|<\alpha^{-p_1+1}I_{p_1-1}(\sigma_n\alpha),\quad y<0,\quad n>N.

\end{gather}



Учитывая формулы \eqref{ab:eq30} и \eqref{ab:eq29}, произведем оценки общих членов ряда $s_1(x,y)$:

\begin{gather}

\label{ab:eq33}

\Bigl|\frac{a_n\alpha^{p_1}}{I_{p_1}(\sigma\alpha)}y^{-p_1}I_{p_1}(\sigma_ny)x\sin2\pi nx\Bigr|<a_n,\quad y>0, \quad n> N;

\\

\label{ab:eq38}

\Bigl|\frac{a_n\alpha^{p_1}}{I_{p_1}(\sigma\alpha)}(-y)^{-p_1}J_{p_1}(-\sigma_ny)x\sin2\pi nx\Bigr|<a_n,\quad y<0,\quad n>N

\end{gather}

и ряда $s_2(x,y)$:

\begin{multline}

\label{ab:eq38}

\Bigl|\Bigl(\frac{2\pi na_n\alpha^{p_1+1}I_{p_1-1}(\sigma_n\alpha)}{\sigma_nI_{p_1}^2(\sigma_n\alpha)}+\frac{b_n\alpha^{p_1}}{I_{p_1}(\sigma_n\alpha)}\Bigr)y^{-p_1}I_{p_1}(\sigma_ny)\cos2\pi nx\Bigr|<\\

<\frac{2\pi n\alpha I_{p_1-1}(\sigma_n\alpha)}{\sigma_n I_{p_1}(\sigma_n\alpha)}+b_n,\quad y>0, \quad n >N;

\end{multline}

\begin{multline}

\label{ab:eq38}

\Bigl|\Bigl(\frac{2\pi na_n\alpha^{p_1+1}I_{p_1-1}(\sigma_n\alpha)}{\sigma_nI_{p_1}^2(\sigma_n\alpha)}+\frac{b_n\alpha^{p_1}}{I_{p_1}(\sigma_n\alpha)}-\frac{2(2p_1+1)\pi n\alpha^{p_1}}{\sigma_n^2I_{p_1}(\sigma_n\alpha)}\Bigr) \times \\

\times (-y)^{-p_1}J_{p_1}(-\sigma_ny)\cos2\pi nx\Bigr|<\\

<\frac{2\pi n\alpha I_{p_1-1}(\sigma_n\alpha)}{\sigma_n I_{p_1}(\sigma_n\alpha)}+b_n+\frac{2p(2p_1+1)\pi na_n}{p_1\sigma_n^2},\quad y<0,\quad n>N .

\end{multline}







Используя неравенства \eqref{ab:eq31} и \eqref{ab:eq32}, получаем оценки общего члена ряда $s_3(x,y)$:

\begin{multline}

\label{ab:eq38}

\Bigl|\frac{2\pi na_n\alpha^{p_1}}{\sigma_nI_{p_1}(\sigma_n\alpha)}y^{-p_1+1}I_{p_1-1}(\sigma_n y)\cos2\pi nx\Bigr| < \\

<\frac{2\pi na_n\alpha I_{p_1-1}(\sigma_n\alpha)}{\sigma_n I_{p_1}(\sigma_n\alpha)},\quad y>0, \quad n >N;

\end{multline}

\begin{multline}

\label{ab:eq38}

\Bigl|\frac{\pi na_n\alpha^{p_1}}{p_1\sigma_nI_{p_1}(\sigma_n\alpha)}(-y)^{-p_1+1}J_{p_1-1}(-\sigma_n y)\cos2\pi nx\Bigr|< \\

<\frac{\pi na_n\alpha I_{p_1-1}(\sigma_n\alpha)}{p_1\sigma_nI_{p_1}(\sigma_n\alpha)},\quad y<0, \quad n>N.

\end{multline}





Учитывая поведение коэффициентов ряда Фурье $m$ раз дифференцируемой функции \cite[c.~636]{ab:bib:Zorich} 

$$

a_n=o\left(n^{-2}\right),\quad b_n=o\left(n^{-2}\right)$$ 

и асимптотику \eqref{ab:eq28}, можно сделать вывод, что ряды, общими членами которых являются правые части неравенств \eqref{ab:eq33}--\eqref{ab:eq38}, сходятся, из чего следует равномерная сходимость рядов $s_i(x,y)$, $i=1, 2, 3,$ на множестве $\overline{D}$.







Докажем выполнение положения \hyperlink{ab:2}{2)}.

Подставим в ряд 

$$

\sum\limits_{n=1}^{\infty}v_n(y)x\sin2\pi nx

$$ 

вместо функций $v_n(y)$  значения, определяемые формулами \eqref{ab:eq13} и \eqref{ab:eq14},  и почленно продифференцируем его  по $y$:

\begin{gather}

s_4(x,y)=\sum\limits_{n=1}^{\infty}\frac{a_n\sigma_n\alpha^{p_1}}{I_{p_1}(\sigma\alpha)}y^{-p_1}I_{p_1+1}(\sigma_ny)x\sin2\pi nx,\quad y>0,

\\

s_4(x,y)=\sum\limits_{n=1}^{\infty}\frac{a_n\sigma_n\alpha^{p_1}}{I_{p_1}(\sigma\alpha)}(-y)^{-p_1}J_{p_1+1}(-\sigma_ny)x\sin2\pi nx,\quad y<0.

\end{gather}



Как было показано в доказательстве  положения \hyperlink{ab:1}{1)},  функция $z^{-\nu}I_{\nu}(z)$ положительна и монотонно возрастает, поэтому в $D_{\varepsilon}$ верна оценка 

\begin{equation}

\label{ab:eq40}

\left|y^{-p_1}I_{p_1+1}(\sigma_ny)\right|<\alpha^{-p_1}I_{p_1+1}(\sigma_n(a-\varepsilon)).

\end{equation}



Так как функция $z^{-\xi} J_{\nu}(z)$ при $\xi\le\nu$ ограничена при $z>0$, справедливо неравенство

\begin{equation}

\label{ab:eq41}

\left|(-y)^{-p_1}J_{p_1+1}(-\sigma_ny)\right|<M\sigma_n^{p_1},

 \end{equation}

  где $$M=\max\limits_{z>0}z^{-p_1}J_{p_1+1}(z).$$

  

Из асимптотики \eqref{ab:eq28} и неравенства \eqref{ab:eq41} получаем, что существует номер $N$ такой, что для $n>N$ выполняется соотношение

\begin{equation}

\label{ab:eq42}

\left|(-y)^{-p_1}J_{p_1+1}(-\sigma_ny)\right|<\alpha^{-p_1}I_{p_1+1}(\sigma_n(\alpha-\varepsilon)).

\end{equation}



Из формул \eqref{ab:eq40} и \eqref{ab:eq42} следует 

\begin{multline}

\label{ab:eq39}

\Bigl|\frac{a_n\sigma_n\alpha^{p_1}}{I_{p_1}(\sigma\alpha)}y^{-p_1}I_{p_1+1}(\sigma_ny)x\sin2\pi nx\Bigr|< \\

<\frac{a_n\sigma_nI_{p_1+1}(\sigma_n(\alpha-\varepsilon))}{I_{p_1}(\sigma\alpha)},\quad y>0, \quad n> N;

\end{multline}

\begin{multline}

\label{ab:eq43}

\Bigl|\frac{a_n\sigma_n\alpha^{p_1}}{I_{p_1}(\sigma\alpha)}(-y)^{-p_1}J_{p_1+1}(-\sigma_ny)x\sin2\pi nx\Bigr|<\\

<\frac{a_n\sigma_nI_{p_1+1}(\sigma_n(\alpha-\varepsilon))}{I_{p_1}(\sigma\alpha)},\quad y<0,\quad n>N.

\end{multline}



В силу асимптотики \eqref{ab:eq28} последовательность 

$$\frac{I_{p_1+1}(\sigma_n(\alpha-\varepsilon))}{I_{p_1}(\sigma\alpha)}$$ 

убывает экспоненциально, остальные множители, входящие в правую часть неравенств \eqref{ab:eq39} и \eqref{ab:eq43}, имеют степенной рост, поэтому ряд, имеющий своим общим членом правую часть неравенств \eqref{ab:eq39} и \eqref{ab:eq43}, сходится, а значит ряд $s_4(x,y)$ сходится равномерно в $D_{\varepsilon}$.



Доказательство равномерной сходимости ряда 

$$\sum\limits_{n=1}^{\infty}v_n(y)x\sin2\pi nx,$$ почленно продифференцированного по $y$ дважды, однократно и двукратно почленно продифференцированного по $x$, а также равномерной сходимости ряда 

$$\sum\limits_{n=1}^{\infty}u_n(y)\cos2\pi nx,$$ 

однократно и двукратно почленно продифференцированного по $x$ и по $y$ соответственно, проводится аналогичным образом. 



Теорема \hyperlink{ab:th2}{2} доказана.

\end{proof}







\thanks{{\bf ORCID}



{\bf Anton Abashkin:} \url{http://orcid.org/0000-0002-3610-1503}

}



\RusRef





\begin{thebibliography}{99}



\RBibitem{ab:bib:Pulkin}

\by Пулькин~С.~П.

\paper О~единственности решения сингулярной задачи Геллерстедта

\jour Изв. вузов. Матем.

\yr 1960

\issue 6

\pages 214--225

\mathnet{http://mi.mathnet.ru/ivm2346}

\mathscinet{http://www.ams.org/mathscinet-getitem?mr=136876}

\zmath{https://zbmath.org/?q=an:0100.09002}



\RBibitem{ab:bib:Volkodavov}

\paper О единственности решения задачи $TN$ для одного уравнения смешанного типа

\by Волкодавов~В.~Ф.

\inbook Волжский математический сборник

\publaddr Куйбышев

\volinfo Вып.~9

\yr 1970

\pages 55--65









\RBibitem{ab:bib:Repin}

\by Репин~О.~А.

\paper Нелокальная задача для уравнения смешанного типа с~сингулярным коэффициентом

\jour  Вестн. Сам. гос. техн. ун-та. Сер. Физ.-мат. науки

\yr 2005

\issue  34

\pages 5--9

\mathnet{http://mi.mathnet.ru/vsgtu331}

\crossref{10.14498/vsgtu331}





\RBibitem{ab:bib:Ruziev}

\by Рузиев~М.~Х.

\paper О нелокальной задаче для уравнения смешанного типа с~сингулярным коэффициентом в~неограниченной области

\jour Изв. вузов. Матем.

\yr 2010

\issue 11

\pages 41--49

\mathnet{http://mi.mathnet.ru/ivm7149}

\mathscinet{http://www.ams.org/mathscinet-getitem?mr=2814563}

%\transl

%\jour Russian Math. (Iz. VUZ)

%\yr 2010

%\vol 54

%\issue 11

%\pages 36--43

%\crossref{http://dx.doi.org/10.3103/S1066369X10110046}



\RBibitem{ab:bib:Sabitov}

\by Сабитов~К.~Б.

\paper Задача Дирихле для уравнений смешанного типа в прямоугольной области

\jour Докл. РАН

\yr 2007

\vol 413

\issue 1

\pages 23--26

\mathnet{http://mi.mathnet.ru/dan701}

\mathscinet{http://www.ams.org/mathscinet-getitem?mr=2447063}



\RBibitem{ab:bib:SabitovSuleimanova}

\by Сабитов~К.~Б., Сулейманова~А.~Х.

\paper Задача Дирихле для уравнения смешанного типа второго рода в~прямоугольной области

\jour Изв. вузов. Матем.

\yr 2007

\issue 4

\pages 45--53

\mathnet{http://mi.mathnet.ru/ivm1382}

\mathscinet{http://www.ams.org/mathscinet-getitem?mr=2380814}

\zmath{https://zbmath.org/?q=an:1158.35391}

%\transl

%\jour Russian Math. (Iz. VUZ)

%\yr 2007

%\vol 51

%\issue 4

%\pages 42--50

%\crossref{http://dx.doi.org/10.3103/S1066369X07040068}





\RBibitem{ab:bib:Il'inMoiseev}

\by Ильин~В.~А., Моисеев~Е.~И.

\paper Двумерная нелокальная краевая задача для оператора Пуассона в~дифференциальной и разностной трактовках

\jour Матем. моделирование

\yr 1990

\vol 2

\issue 8

\pages 139--156

\mathnet{http://mi.mathnet.ru/mm2433}

\mathscinet{http://www.ams.org/mathscinet-getitem?mr=1086120}

\zmath{https://zbmath.org/?q=an:0993.35500}





\RBibitem{ab:bib:Francl}

\paper К образованию скачков уплотнения в дозвуковых течениях с местными сверхзвуковыми скоростями

\by Франкль~Ф.~И.

\jour Прикл. матем. и мех.

\yr 1947

\vol 11

\issue 1

\pages 199--202

\zmath{https://zbmath.org/?q=an:0041.54305}



\RBibitem{ab:bib:Moiseev2}

\paper О решении спектральным методом одной нелокальной краевой задачи

\by Моисеев~Е.~И.

\jour Диффер. уравн.

\yr 1999

\vol 35

\issue 8

\pages 1094--1100

\zmath{https://zbmath.org/?q=an:0973.35085}

%Moiseev, E.I.

%The solution of a nonlocal boundary value problem by the spectral method. (English. Russian original) Zbl 0973.35085

%Differ. Equations 35, No.8, 1105-1112 (1999); translation from Differ. Uravn. 35, No.8, 1094-1100 (1999).





\RBibitem{ab:bib:LernerRepin2}

\paper О задачах типа задачи Франкля для некоторых эллиптических уравнений с вырождением разного рода

\by Лернер~М.~Е., Репин~О.~А.

\jour Диффер. уравн.

\yr 1999

\vol 35

\issue 8

\pages 1087--1093

\zmath{https://zbmath.org/?q=an:0976.35028}

%Lerner, M.E.; Repin, O.A.

%On Frankl type problems for some elliptic equations with degeneration of various kinds. (English. Russian original) Zbl 0976.35028

%Differ. Equations 35, No.8, 1098-1104 (1999); translation from Differ. Uravn. 35, No.8, 1087-1093 (1999).





\RBibitem{ab:bib:SabitovSidorenko}

\by Сабитов~К.~Б., Сидоренко~О.~Г.

\paper Нелокальная задача для вырождающегося гиперболического уравнения

\inbook Современные проблемы физики и математики

\bookvol 1

\bookinfo Труды Всероссийской научной конференции 

\procinfo 16--18 сентября 2004~г.

\publaddr Уфа

\publ Гилем

\yr 2004

\pages 80--86





\RBibitem{ab:bib:Sidorenko}

\by Сидоренко~О.~Г.

\paper Существенно нелокальная задача для уравнения смешанного типа в~полуполосе

\jour Изв. вузов. Матем.

\yr 2007

\issue 3

\pages 60--64

\mathnet{http://mi.mathnet.ru/ivm1373}

\mathscinet{http://www.ams.org/mathscinet-getitem?mr=2336877}

\zmath{https://zbmath.org/?q=an:1142.35051}

%\transl

%\jour Russian Math. (Iz. VUZ)

%\yr 2007

%\vol 51

%\issue 3

%\pages 55--59

%\crossref{http://dx.doi.org/10.3103/S1066369X07030085}





\RBibitem{ab:bib:LernerRepin1}

\paper Нелокальные краевые задачи в вертикальной полуполосе для обобщенного осесимметрического уравнения Гельмгольца

\by Лернер~М.~Е., Репин~О.~А.

\jour Диффер. уравн.

\yr 2001

\vol 37

\issue 11

\pages 1562--1564

\zmath{https://zbmath.org/?q=an:1016.35016}

%Lerner, M.E.; Repin, O.A.

%Nonlocal boundary value problems in a vertical half-strip for a generalized axisymmetric Helmholtz equation. (English. Russian original) Zbl 1016.35016

%Differ. Equ. 37, No.11, 1640-1642 (2001); translation from Differ. Uravn. 37, No.11, 1562-1564 (2001).

%\crossref{10.1023/A:1017985319783}



\RBibitem{ab:bib:Moiseev1}

\paper О разрешимости одной нелокальной краевой задачи

\by Моисеев~Е.~И.

\jour Диффер. уравн.

\yr 2001

\issue 11

\vol 37

\pages 1565--1567

%Moiseev, E.I.

%Solvability of a nonlocal boundary value problem. (English. Russian original) Zbl 1016.35021

%Differ. Equ. 37, No.11, 1643-1646 (2001); translation from Differ. Uravn. 37, No.11, 1565-1567 (2001).

%\crossref{10.1023/A:1017937403853}



\RBibitem{ab:bib:Abashkin1}

\by Абашкин~А.~А.

\paper Однозначная разрешимость нелокальной задачи для осесимметрического уравнения Гельмгольца

\jour Вестн. СамГУ. Естественнонаучн. сер.

\yr 2011

\issue 2(83)

\pages 5--14

\mathnet{http://mi.mathnet.ru/vsgu43}



\RBibitem{ab:bib:Abashkin2}

\by Абашкин~А.~А.

\paper Об одной нелокальной задаче для осесимметрического уравнения Гельмгольца

\jour Вестн. Сам. гос. техн. ун-та. Сер. Физ.-мат. науки

\yr 2011

\issue 3(24)

\pages 26--34

\mathnet{http://mi.mathnet.ru/vsgtu852}

\crossref{10.14498/vsgtu852}







\RBibitem{ab:bib:Sabitova1}

\paper Нелокальная задача для уравнения Лаврентьева--Бицадзе в прямоугольной области

\by Сабитова~Ю.~К.

\inbook Труды Стерлитамакского филиала АН РБ

\volinfo Вып.~6

\publaddr Уфа

\publ Гилем

\yr 2009

\pages 94--102

%Yu. K. Sabitova and A. A. Bakhristova, "The Dirichlet Problem for a Mixed Type Equation with the Lavrent'ev-Bitsadze Operator," Trudy Sterlitamakskogo Filiala AN RB, Ser. Phys.-Mat. and Tekhn. Nauki, No. 6, 103-110 (Gilem, Ufa, 2009).



\RBibitem{ab:bib:Sabitova2}

\by Сабитова~Ю.~К.

\thesis Краевые задачи с нелокальным условием для уравнений смешанного типа в прямоугольной области 

\thesisinfo Дисс. \dots\ к.ф.м.н

\publaddr Стерлитамак

\yr 2007

\totalpages 133

\elib{http://elibrary.ru/item.asp?id=16110932}

% Yu. K. Sabitova, “Boundary-value problems with nonlocal condition for mixed-type equations in a rectangular domain,” Theses of Candidate Dissertation in Physics and Mathematics, SGPA, Sterlitamak (2007).





\RBibitem{ab:bib:Lebedev}

\by Лебедев~Н.~Н.

\book Специальные функции и их приложения

\publaddr Спб.

\publ Лань

\yr 2010

\totalpages 368

%Lebedev, N.N.

%Special functions and their applications. Rev. engl. ed. Translated and edited by Richard A. Silverman. (English) Zbl 0271.33001

%New York: Dover Publications, Inc. XI, 308 p. $ 3.75 (1972).

%\zmath{https://zbmath.org/?q=an:0271.33001}





\RBibitem{ab:bib:Zorich}

\by Зорич~В.~А.

\book Математический анализ

\bookvol 2

\publaddr М.

\publ МЦНМО

\yr 2002

\totalpages 787

%Zorich, Vladimir A.

%Mathematical analysis II. Transl. from the 4th Russian edition by Roger Cooke. (English) Zbl 1071.00003

%Universitext. Berlin: Springer (ISBN 3-540-40633-6/hbk). xv, 681 p. EUR 49.95/net; sFr. 86.00; ? 38.50; $ 69.95 (2004).

%\zmath{https://zbmath.org/?q=an:1071.00003}





\end{thebibliography}









\RuDate	



\bigskip





\makeentitle



\thanks{{\bf ORCID}





{\bf Anton Abashkin:} \url{http://orcid.org/0000-0002-3610-1503}

}







\EngRef















\begin{thebibliography}{99}



\Bibitem{ab:bib:Pulkin-eng}

\lang In Russian

\by Pul'kin~S.~P.

\paper Uniqueness of the solution of a~singular problem of Gellerstedt--Tricomi

\jour Izv. Vyssh. Uchebn. Zaved. Mat.

\yr 1960

\issue 6

\pages 214--225





\Bibitem{ab:bib:Volkodavov-eng}

\lang In Russian

\paper The uniqueness of the solution of the $TN$ problem for a certain equation of mixed type

\by Volkodavov~V.~F.

\inbook Volzhskii matematicheskii sbornik

\publaddr Kuibyshev

\volinfo Vyp.~9

\yr 1970

\pages 55--65







\Bibitem{ab:bib:Repin-eng}

\lang In Russian

\by Repin~O.~A.

\paper A nonlocal problem for a mixed-type equation with a singular coefficient

\jour  Vestn. Samar. Gos. Tekhn. Univ. Ser. Fiz.-Mat. Nauki

\yr 2005

\issue 34

\pages 5--9

\crossref{10.14498/vsgtu331}





\Bibitem{ab:bib:Ruziev-eng}

\paper A nonlocal problem for a mixed-type equation with a singular coefficient in an unbounded domain

\by Ruziev~M.~Kh.

\mathscinet{http://www.ams.org/mathscinet-getitem?mr=2814563}

\jour Russian Math. (Iz. VUZ)

\yr 2010

\vol 54

\issue 11

\pages 36--43

\crossref{10.3103/S1066369X10110046}

\scopus{2-s2.0-78649583445}



\Bibitem{ab:bib:Sabitov-eng}

\lang In Russian

\by Sabitov~K.~B.

\paper The Dirichlet problem for equations of mixed type in a rectangular domain

\jour Dokl. Akad. Nauk 

\vol 413 

\yr 2007

\issue 1

\pages 23--26





\Bibitem{ab:bib:SabitovSuleimanova-eng}

\paper The Dirichlet problem for a mixed-type equation of the second kind in a rectangular domain

\by Sabitov~K.~B., Suleimanova~A.~Kh.

\mathscinet{http://www.ams.org/mathscinet-getitem?mr=2380814}

\zmath{https://zbmath.org/?q=an:1158.35391}

\jour Russian Math. (Iz. VUZ)

\yr 2007

\vol 51

\issue 4

\pages 42--50

\crossref{10.3103/S1066369X07040068}





\Bibitem{ab:bib:Il'inMoiseev-eng}

\lang In Russian

\by Il'in~V.~A., Moiseev~E.~I.

\paper A two-dimensional nonlocal boundary value problem for the Poisson operator in the differential and the difference interpretation

\jour Matem. Modelirovanie

\yr 1990

\vol 2

\issue 8

\pages 139--156

\mathscinet{http://www.ams.org/mathscinet-getitem?mr=1086120}

\zmath{https://zbmath.org/?q=an:0993.35500}





\Bibitem{ab:bib:Francl-eng}

\lang In Russian

\by Frankl~F.~I.

\paper Causes of shock waves in subsonic flows with local supersonic velocities

\jour Prikl. Matem. Mekh.

\yr 1947

\vol 11

\issue 1

\pages 199--202









\Bibitem{ab:bib:Moiseev2-eng}

\zmath{https://zbmath.org/?q=an:0973.35085}

\by Moiseev~E.~I.

\paper The solution of a nonlocal boundary value problem by the spectral method

\jour Differ. Equ.

\vol 35

\issue 8

\pages 1105--1112 

\yr 1999

\scopus{2-s2.0-27144519300}





\Bibitem{ab:bib:LernerRepin2-eng}

\zmath{https://zbmath.org/?q=an:0976.35028}

\by Lerner~M.~E.,  Repin~O.~A.

\paper On Frankl type problems for some elliptic equations with degeneration of various kinds

\jour Differ. Equ.

\vol 35

\issue 8

\pages 1098--1104 

\yr 1999

\scopus{2-s2.0-27144434883}





\Bibitem{ab:bib:SabitovSidorenko-eng}

\publaddr Ufa

\publ Gilem

\yr 2004

\pages 80--86

\bookvol 1

\inbook Sovremennye problemy fiziki i matematiki \rm [Contemporary Problems in Physics and Mathematics]

\by Sabitov~K.~B., Sidorenko~O.~G.

\paper Solvability of a Nonlocal Boundary-Value Problem

\bookinfo Proceedings of All-Russia Conference 

\lang In Russian







\Bibitem{ab:bib:Sidorenko-eng}

\paper An essentially nonlocal problem for a mixed-type equation in a semiband

\by  Sidorenko~O.~G.

\jour Russian Math. (Iz. VUZ)

\yr 2007

\vol 51

\issue 3

\pages 55--59

\crossref{10.3103/S1066369X07030085}





\Bibitem{ab:bib:LernerRepin1-eng}

\zmath{https://zbmath.org/?q=an:1016.35016}

\by Lerner~M.~E., Repin~O.~A.

\paper Nonlocal boundary value problems in a vertical half-strip for a generalized axisymmetric Helmholtz equation

\jour Differ. Equ. 

\vol 37

\issue 11

\pages 1640--1642 

\yr 2001

\crossref{10.1023/A:1017985319783}

\scopus{2-s2.0-0035567008}



\Bibitem{ab:bib:Moiseev1-eng}

\by Moiseev~E.~I.

\paper Solvability of a nonlocal boundary value problem

\jour Differ. Equ.

\vol 37

\issue 11

\pages 1643--1646 

\yr 2001

\crossref{10.1023/A:1017937403853}

\scopus{2-s2.0-0035564095}



\Bibitem{ab:bib:Abashkin1-eng}

\by Abashkin~A.~A.

\paper One-valued solvability of a~nonlocal problem for the axisymmetric Helmholtz equation

\jour Vestnik SamGU. Estestvenno-Nauchnaya Ser.

\yr 2011

\issue 2(83)

\pages 5--14

\lang In Russian



\Bibitem{ab:bib:Abashkin2-eng}

\lang In Russian

\by Abashkin~A.~A.

\paper On one non-local problem for axisymmetric Helmholtz equation

\jour Vestn. Samar. Gos. Tekhn. Univ. Ser. Fiz.-Mat. Nauki

\yr 2011

\issue 3(24)

\pages 26--34

\crossref{10.14498/vsgtu852}









\Bibitem{ab:bib:Sabitova1-eng}

\lang In Russian

\by Sabitova~Yu.~K.

\paper Nonlocal problem for Lavrent'ev--Bitsadze equation in a rectangular region

\inbook Trudy Sterlitamakskogo filiala AN RB

\volinfo no.~6

\publaddr Ufa

\publ Gilem

\yr 2009

\pages 94--102







\Bibitem{ab:bib:Sabitova2-eng}

\lang In Russian

\by Sabitova~Yu.~K.

\thesis Boundary-value problems with nonlocal condition for mixed-type equations in a rectangular domain

\thesisinfo Theses of Candidate Dissertation in Physics and Mathematics

\publaddr Sterlitamak 

\yr 2007

\totalpages 133









\Bibitem{ab:bib:Lebedev-eng}

\by Lebedev~N.~N.

\book Special functions and their applications

\publaddr New York

\publ Dover Publications, Inc

\totalpages xi+308 

\yr 1972

\zmath{https://zbmath.org/?q=an:0271.33001}





\Bibitem{ab:bib:Zorich-eng}

\by Zorich~V.~A.

\book Mathematical analysis 

\bookvol II

\serial Universitext

\publaddr Berlin

\yr 2004

\publ Springer 

\totalpages xv+681 

\zmath{https://zbmath.org/?q=an:1071.00003}





\end{thebibliography}



\EngDate





\newpage\Russian



%

%

%\end{document}

